\section{AHSME 1981}
        \begin{enumerate}
        	\item (\textit{AHSME 1981}) If $\sqrt{x+2}=2$, then $(x+2)^{2}$ equals


 $\textbf{(A)}\ \sqrt{2}\qquad\textbf{(B)}\ 2\qquad\textbf{(C)}\ 4\qquad\textbf{(D)}\ 8\qquad\textbf{(E)}\ 16$


 
 \textbf{Answer: }E



\textbf{Solution 1}

Square both sides of the equation to get get $x+2 = 4$.  Then square both sides again to get $(x+2)^2 = \boxed{\textbf{(E) }16}$

 



\textbf{Solution 2}

Solve the equation to find that $x = 2$. Then $(x+2)^2=(2+2)^2=4^2=\boxed{\textbf{(E) }16}$

 


	\item (\textit{AHSME 1981}) Point $E$ is on side $AB$ of square $ABCD$. If $EB$ has length one and $EC$ has length two, then the area of the square is


 
                    \begin{figure}[H]
                        centering
                        \includegraphics[width=0.5\textwidth]{Problem/AHSME_I_1981_2_Problem_Figure_1.png}
                    \end{figure}
                    


 

$\textbf{(A)}\ \sqrt{3}\qquad\textbf{(B)}\ \sqrt{5}\qquad\textbf{(C)}\ 3\qquad\textbf{(D)}\ 2\sqrt{3}\qquad\textbf{(E)}\ 5$


 
 \textbf{Answer: }C



\textbf{Solution}

Note that $\triangle BCE$ is a right triangle. By the Pythagorean theorem, $BC^2 = CE^2 - BE^2 = 2^2-1^2=3$, so the area of $ABCD$ is  $\boxed{\textbf{(C)}\ 3}$. 

 


	\item (\textit{AHSME 1981}) For $x\neq0$, $\dfrac{1}{x}+\dfrac{1}{2x}+\dfrac{1}{3x}$ equals


 $\textbf{(A)}\ \dfrac{1}{2x}\qquad\textbf{(B)}\ \dfrac{1}{6x}\qquad\textbf{(C)}\ \dfrac{5}{6x}\qquad\textbf{(D)}\ \dfrac{11}{6x}\qquad\textbf{(E)}\ \dfrac{1}{6x^3}$


 
 \textbf{Answer: }D



\textbf{Solution}

The least common multiple of $x$, $2x$, and $3x$ is $6x$, so $\frac{1}{x} + \frac{1}{2x} + \frac{1}{3x} = \frac{6}{6x}$ + $\frac{3}{6x}$ + $\frac{2}{6x} = \boxed{\left(\mathbf{D}\right) \frac{11}{6x}}$.

 


	\item (\textit{AHSME 1981}) If three times the larger of two numbers is four times the smaller and the difference between the numbers is $8$, the the larger of two numbers is 


 $\textbf{(A)}\ 16\qquad\textbf{(B)}\ 24\qquad\textbf{(C)}\ 32\qquad\textbf{(D)}\ 44\qquad\textbf{(E)}\ 52$


 
 \textbf{Answer: }C



\textbf{Solution}

Let $x$ be the smaller number and $y$ be the larger number.   Then $3y = 4x$, so $x = \frac{3}{4}y$.  Also, $y - x = 8$, so $y - \frac{3}{4}y = \frac{1}{4}y = 8$, and $y = \boxed{(\textbf{C})\ 32}$.

 


	\item (\textit{AHSME 1981}) In trapezoid $ABCD$, sides $AB$ and $CD$ are parallel, and diagonal $BD$ and side $AD$ have equal length. If $\angle DCB=110^\circ$ and $\angle CBD=30^\circ$, then $\angle ADB=$


 
                    \begin{figure}[H]
                        centering
                        \includegraphics[width=0.5\textwidth]{Problem/AHSME_I_1981_5_Problem_Figure_1.png}
                    \end{figure}
                    


 $\textbf{(A)}\ 80^\circ\qquad\textbf{(B)}\ 90^\circ\qquad\textbf{(C)}\ 100^\circ\qquad\textbf{(D)}\ 110^\circ\qquad\textbf{(E)}\ 120^\circ$


 
 \textbf{Answer: }C



\textbf{Solution}

In $\triangle BCD$, $\angle CDB=180^\circ - \angle DCB - \angle CBD = 180^\circ - 30^\circ - 110^\circ = 40^\circ.$

 Because $AB \parallel CD$, $\angle CDB= \angle DBA = 40^\circ$. Since $\triangle DAB$ is isosceles, $\angle ADB = 180^\circ - 2 \cdot 40^\circ = \boxed{(\textbf{C})\ 100^\circ}$.

 -edited by coolmath34, j314andrews

 


	\item (\textit{AHSME 1981}) If $\dfrac{x}{x-1}=\dfrac{y^2+2y-1}{y^2+2y-2}$, then $x$ equals 


 $\textbf{(A)}\ y^2+2y-1\qquad\textbf{(B)}\ y^2+2y-2\qquad\textbf{(C)}\ y^2+2y+2 \qquad \\ \textbf{(D)}\ y^2+2y+1\qquad\textbf{(E)}\ -y^2-2y+1$


 
 \textbf{Answer: }A



\textbf{Solution}

Cross-multiply and simplify:
 \[x(y^2 + 2y - 2) = (x-1)(y^2 + 2y - 1)\]
 \[xy^2 + 2xy - 2x = xy^2 + 2xy - x - y^2 - 2y + 1\]
 \[-x = -y^2 - 2y + 1\]
 \[x = \boxed{(\textbf{A})\ y^2 + 2y - 1}\]

 -edited by coolmath34

 


 


 



\textbf{Solution 2}


 \[\frac{x}{x-1} = \frac{y^2+2y-1}{y^2+2y-2} \Rightarrow \frac{x+(x-1)}{x-(x-1)} = \frac{(y^2+2y-1)+(y^2+2y-2)}{(y^2+2y-1)-(y^2+2y-2)}\]

 


	\item (\textit{AHSME 1981}) How many of the first one hundred positive integers are divisible by all of the numbers $2$, $3$, $4$, and $5$?


 $\textbf{(A)}\ 0\qquad\textbf{(B)}\ 1\qquad\textbf{(C)}\ 2\qquad\textbf{(D)}\ 3\qquad\textbf{(E)}\ 4$


 
 \textbf{Answer: }B



\textbf{Solution}

The least common multiple of $2$, $3$, $4$, and $5$ is $2^2 \cdot 3 \cdot 5 = 60$. There is only one multiple of $60$ between $1$ and $100$, so the answer is $\boxed{\textbf{(B)}\ 1}$.

 -edited by coolmath34

 


	\item (\textit{AHSME 1981}) For all positive numbers $x$, $y$, $z$, the product $(x+y+z)^{-1}(x^{-1}+y^{-1}+z^{-1})(xy+yz+xz)^{-1}[(xy)^{-1}+(yz)^{-1}+(xz)^{-1}]$ equals $\textbf{(A)}\ x^{-2}y^{-2}z^{-2}\qquad\textbf{(B)}\ x^{-2}+y^{-2}+z^{-2}\qquad\textbf{(C)}\ (x+y+z)^{-1}\qquad \textbf{(D)}\ \dfrac{1}{xyz}\qquad \\ \textbf{(E)}\ \dfrac{1}{xy+yz+xz}$


 
 \textbf{Answer: }A



\textbf{Solution}

Converting all negative exponents to fractions yields:
 \[\left(\frac{1}{x+y+z}\right)\left(\frac{1}{x}+\frac{1}{y}+\frac{1}{z}\right)\left(\frac{1}{xy+yz+xz}\right)\left(\frac{1}{xy}+\frac{1}{xz}+\frac{1}{yz}\right)\]Adding fractions within each factor yields:
 \[\left(\frac{1}{x+y+z}\right)\left(\frac{xy+yz+xz}{xyz}\right)\left(\frac{1}{xy+yz+xz}\right)\left(\frac{x+y+z}{xyz}\right)\]Multiplying and cancelling common factors yields:
 \[\frac{1}{x^2 y^2 z^2} = \boxed{\textbf{(A)}\ x^{-2}y^{-2}z^{-2}}\]

 
-edited by coolmath34

 


	\item (\textit{AHSME 1981}) In the adjoining figure, $PQ$ is a diagonal of the cube. If $PQ$ has length $a$, then the surface area of the cube is


 
                    \begin{figure}[H]
                        centering
                        \includegraphics[width=0.5\textwidth]{Problem/AHSME_I_1981_9_Problem_Figure_1.png}
                    \end{figure}
                    


 $\textbf{(A)}\ 2a^2\qquad\textbf{(B)}\ 2\sqrt{2}a^2\qquad\textbf{(C)}\ 2\sqrt{3}a^2\qquad\textbf{(D)}\ 3\sqrt{3}a^2\qquad\textbf{(E)}\ 6a^2$


 
 \textbf{Answer: }A



\textbf{Solution}

Let $s$ be the side length of the cube.  By the Pythagorean theorem, any diagonal of any face of the cube has length $s\sqrt{2}$, and thus any diagonal of the cube has length $a = s \sqrt{3}$. So $s = \frac{a}{\sqrt{3}}$, and therefore the surface area of the cube is $6\left(\frac{a}{\sqrt{3}}\right)^2=6 \cdot \frac{a^2}{3}=\boxed{(\textbf{A})\ 2a^2}$.

 


	\item (\textit{AHSME 1981}) The lines $L$ and $K$ are symmetric to each other with respect to the line $y=x$. If the equation of the line $L$ is $y=ax+b$ with $a\neq 0$ and $b\neq 0$, then the equation of $K$ is $y=$


 $\textbf{(A)}\ \dfrac{1}{a}x+b\qquad\textbf{(B)}\ -\dfrac{1}{a}x+b\qquad\textbf{(C)}\ \dfrac{1}{a}x-\dfrac{b}{a}\qquad\textbf{(D)}\ \dfrac{1}{a}x+\dfrac{b}{a}\qquad\textbf{(E)}\ \dfrac{1}{a}x-\dfrac{b}{a}$


 
 \textbf{Answer: }E



\textbf{Solution}

Recall that the reflection of any point $(p, q)$ across the line $y = x$ is $(q, p)$.  Therefore, any point $(p, q)$ lies on $L$ if and only if $(q, p)$ lies on $K$. Therefore, any point $(r, s)$ lies on $K$ if and only if $(s, r)$ lies on $L$, that is, $r=as+b$.  Therefore, line $K$ has equation $x = ay + b$, and isolating $y$ yields $y = \boxed{(\textbf{E})\ \frac{1}{a}x-\frac{b}{a}}$.

 


	\item (\textit{AHSME 1981}) The three sides of a right triangle have integral lengths which form an arithmetic progression. One of the sides could have length


 $\textbf{(A)}\ 22\qquad\textbf{(B)}\ 58\qquad\textbf{(C)}\ 81\qquad\textbf{(D)}\ 91\qquad\textbf{(E)}\ 361$


 
 \textbf{Answer: }C



\textbf{Solution}

Let the three sides be $a-d$, $a$, and $a+d$.  By the Pythagorean Theorem,
 \[(a-d)^2 + a^2 = (a+d)^2\]This can be simplified to
 \[a(a-4d) = 0\]which has solutions
 \[a = 0, a = 4d\]

 Since $a=0$ is invalid, $a=4d$ and the triangle has sides $3d$, $4d$, $5d$.  Therefore, the correct answer must be divisible by $3$, $4$, or $5$. The only valid answer choice is $\boxed{\textbf{(C)}\ 81}$, since it is divisible by $3$.

 -edited by coolmath34

 


	\item (\textit{AHSME 1981}) If $p$, $q$, and $M$ are positive numbers and $q<100$, then the number obtained by increasing $M$ by $p\%$ and decreasing the result by $q\%$ exceeds $M$ if and only if


 $\textbf{(A)}\ p>q \qquad\textbf{(B)}\ p>\dfrac{q}{100-q}\qquad\textbf{(C)}\ p>\dfrac{q}{1-q}\qquad \textbf{(D)}\ p>\dfrac{100q}{100+q}\qquad\textbf{(E)}\ p>\dfrac{100q}{100-q}$


 
 \textbf{Answer: }E



\textbf{Solution 1}

The goal is to find when $M\left(1+\frac{p}{100}\right)\left(1-\frac{q}{100}\right) > M$.  Isolate $p$ as follows:

 
 \[\left(1+\frac{p}{100}\right)\left(1-\frac{q}{100}\right) > 1\]

 
 \[\left(\frac{100+p}{100}\right)\left(\frac{100-q}{100}\right) > 1\]

 
 \[\frac{(100+p)(100-q)}{10000} > 1\]

 
 \[100+p > \frac{10000}{100-q}\]

 
 \[p > \frac{10000}{100-q} - 100 = \boxed{(\textbf{E})\ \frac{100q}{100-q}}\]

 -j314andrews

 



\textbf{Solution 2 (Answer Choices)}

Choice A is incorrect. If $(M, p, q) = (100, 50, 40)$, then $p > q$, but $M$ will not increase, as $100 \cdot 1.5 \cdot 0.6 = 90$.

 Choice B is incorrect. If $(M, p, q) = (100, 50, 50)$, then $p > \frac{q}{100-q}$, but $M$ will not increase, as $100 \cdot 1.5 \cdot 0.5 = 75$.

 Choice C is incorrect. If $(M, p, q) = (100, 50, 50)$, then $p > \frac{q}{1-q}$, but $M$ will not increase, as $100 \cdot 1.5 \cdot 0.5 = 75$. 

 Choice D is incorrect. If $(M, p, q) = (100, 100, 60)$, then $p > \frac{100q}{100+q}$, but $M$ will not increase, as $100 \cdot 2 \cdot 0.4 = 80$.

 By process of elimination, $\boxed{(\textbf{E})\ \dfrac{100q}{100-q}}$ must be correct.

 -Arcticturn, edited by j314andrews

 


	\item (\textit{AHSME 1981}) Suppose that at the end of any year, a unit of money has lost $10\%$ of the value it had at the beginning of that year. Find the smallest integer $n$ such that after $n$ years, the money will have lost at least $90\%$ of its value (To the nearest thousandth $\log_{10} 3 = 0.477$).


 $\textbf{(A)}\ 14\qquad\textbf{(B)}\ 16\qquad\textbf{(C)}\ 18\qquad\textbf{(D)}\ 20\qquad\textbf{(E)}\ 22$


 
 \textbf{Answer: }E



\textbf{Solution}

After $n$ years, the money will be worth $(0.9)^n$ times what it is currently worth.  Therefore, the goal is to find the smallest integer such that $(0.9)^n < 0.1$.  Taking the base-$10$ logarithm of both sides yields $n \log_{10}{0.9} < \log_{10}{0.1}$, which is equivalent to $n > \frac{\log_{10}{0.1}}{\log_{10}{0.9}} = \frac{-1}{-1+2\log_{10}{3}}$. 

 Substituting $\log_{10}{3} \approx 0.477$ yields $n > \frac{-1}{-1 + 2 \cdot 0.477} = \frac{-1}{-0.046} \approx 21.7$, so the minimum possible integer value of $n$ is $\boxed{(\textbf{E})\ 22}$.

 -edited by Maxxie, maxamc, j314andrews

 


	\item (\textit{AHSME 1981}) In a geometric sequence of real numbers, the sum of the first $2$ terms is $7$, and the sum of the first $6$ terms is $91$. The sum of the first $4$ terms is


 $\textbf{(A)}\ 28\qquad\textbf{(B)}\ 32\qquad\textbf{(C)}\ 35\qquad\textbf{(D)}\ 49\qquad\textbf{(E)}\ 84$


 
 \textbf{Answer: }A



\textbf{Solution 1}

Let $a$ be the first term of the sequence and $r$ be the common ratio.  Then $a+ar = a(1+r) = 7$ and $a+ar+ar^2+ar^3+ar^4+ar^5 = a(1+r)(1+r^2+r^4) = 7(1+r^2+r^4) = 91$.  So $1+r^2+r^4=13$, and $r^4+r^2-12=(r^2-3)(r^2+4)=0$.  Since $r$ is real, $r^2 = 3$, so the sum of the first $4$ terms is $a+ar+ar^2+ar^3 = a(1+r)(1+r^2) = 7\cdot 4 = \boxed{(\textbf{A})\ 28}$.

 



\textbf{Solution 2}

Denote the sum of the first $2$ terms as $x$. Since we know that the sum of the first $6$ terms is $91$ which is $7 \cdot 13$, we have $x$ + $xy$ + $xy^2$ = $13x$ because it is a geometric series. We can quickly see that $y$ = $3$, and therefore, the sum of the first $4$ terms is $4x = 4 \cdot 7 = \boxed {(\textbf{A})\ 28}$

 


	\item (\textit{AHSME 1981}) If $b>1$, $x>0$, and $(2x)^{\log_b 2}-(3x)^{\log_b 3}=0$, then $x$ is


 $\textbf{(A)}\ \dfrac{1}{216}\qquad\textbf{(B)}\ \dfrac{1}{6}\qquad\textbf{(C)}\ 1\qquad\textbf{(D)}\ 6\qquad\textbf{(E)}\ \text{not uniquely determined}$


 
 \textbf{Answer: }B



\textbf{Solution}

Use rules of logarithms to solve this equation.  

 
 \[(2x)^{\log_{b} 2} = (3x)^{\log_{b} 3}\] 

 
 \[\log_{b} 2 \cdot \log_{b} 2x = \log_{b} 3 \cdot \log_{b} 3x\] 

 
 \[\log_{b} 2 \cdot (\log_{b} 2 + \log_{b} x) = \log_{b} 3 \cdot (\log_{b} 3 + \log_{b} x)\]

 
 \[(\log_{b} 2)^2  + \log_{b} 2 \cdot \log_{b} x = (\log_{b} 3)^2 + \log_{b} 3 \cdot \log_{b} x\]

 
 \[(\log_{b} 2)^2 - (\log_{b} 3)^2  =  \log_{b} 3 \cdot \log_{b} x - \log_{b} 2 \cdot \log_{b} x\]

 
 \[(\log_{b} 2 + \log_{b} 3)(\log_{b} 2 - \log_{b} 3) = \log_{b} x \cdot (\log_{b} 3 - \log_{b} 2)\]

 
 \[\log_{b} 2 + \log_{b} 3 = -\log_{b} x\]

 
 \[\log_{b} 6 = -\log_{b} x\]

 
 \[6 = \frac{1}{x}\]

 
 \[x = \boxed {(\textbf{B})\ \frac{1}{6}}\]

 -j314andrews

 


	\item (\textit{AHSME 1981}) The base three representation of $x$ is 

 \[12112211122211112222\]
The first digit (on the left) of the base nine representation of $x$ is


 $\textbf{(A)}\ 1\qquad\textbf{(B)}\ 2\qquad\textbf{(C)}\ 3\qquad\textbf{(D)}\ 4\qquad\textbf{(E)}\ 5$


 
 \textbf{Answer: }E



\textbf{Solution 1}

Since $9=3^2$, every $2$ digits in base $3$ corresponds to $1$ digit in base $9$. Since this number is $20$ digits long, which is an even number of digits, the answer must correspond to the first $2$ digits on the left. So the answer is $12_{3} = 1 \cdot 3 + 2 = 5\ \fbox{(E)}$.

 -edited by coolmath34, j314andrews

 



\textbf{Solution 2 (Long Way)}

Convert $x$ to base 10 then convert the result to base 9.
 \[12112211122211112222_{3} = 2150029898\]

 
 \[2150029898 = 5484584488_{9}\]

 Therefore, the answer is $\textbf{(E)}\ 5.$

 -edited by coolmath34

 


	\item (\textit{AHSME 1981}) The function $f$ is not defined for $x=0$, but, for all non-zero real numbers $x$, $f(x)+2f\left(\dfrac{1}x\right)=3x$. The equation $f(x)=f(-x)$ is satisfied by


 $\textbf{(A)}\ \text{exactly one real number} \qquad \textbf{(B)}\ \text{exactly two real numbers} \qquad\textbf{(C)}\ \text{no real numbers}\qquad \\ \textbf{(D)}\ \text{infinitely many, but not all, non-zero real numbers} \qquad\textbf{(E)}\ \text{all non-zero real numbers}$


 
 \textbf{Answer: }B



\textbf{Solution}

Substitute $x$ with $\frac{1}{x}$:

 $f\left(\frac{1}{x}\right)+2f(x)=\frac{3}{x}$.

 Adding this to $f(x)+2f\left(\dfrac{1}x\right)=3x$, we get

 $3f(x)+3f\left(\dfrac{1}x\right)=3x+\frac{3}{x}$, or

 $f(x)+f\left(\dfrac{1}x\right)=x+\frac{1}{x}$.

 Subtracting this from $f(\frac{1}{x})+2f(x)=\frac{3}{x}$, we have

 $f(x)=\frac{2}{x}-x$.

 Then, $f(x)=f(-x)$ when $\frac{2}{x}-x=\frac{2}{-x}+x$ or

 $\frac{2}{x}-x=0$, so $x=\pm{\sqrt2}$ are the two real solutions and the answer is $\boxed{(B)}$.

 


	\item (\textit{AHSME 1981}) The number of real solutions to the equation

 \[\dfrac{x}{100}=\sin x\]
is


 $\textbf{(A)}\ 61\qquad\textbf{(B)}\ 62\qquad\textbf{(C)}\ 63\qquad\textbf{(D)}\ 64\qquad\textbf{(E)}\ 65$


 
 \textbf{Answer: }C



\textbf{Solution}

The answer to this problem is the number of intersections between the graph of $f(x) = \sin x$ and $f(x) = \frac{1}{100}x.$ We can do the right side of the coordinate plane first. Each cycle of the sine wave, consisting of 2π, will have 2 intersections (From the positive part of the sine wave) The line $f(x) = \frac{1}{100}x$ will consist of 16 cycles plus a little bit extra for $x$ from 1 to 100. However, the extra is not complete enough to have any intersection at all. Thus, the number of intersections is $2 \cdot 16 = 32.$ Because of symmetry, we can multiply by two to account for the left side, and subtract one because of the origin. So the answer is $32 \cdot 2 - 1 =  \textbf{(C)}\ 63.$

 \href{https://www.desmos.com/calculator/z6edqwu1kx}{https://www.desmos.com/calculator/z6edqwu1kx} - Graph

 


	\item (\textit{AHSME 1981}) In $\triangle ABC$, $M$ is the midpoint of side $BC$, $AN$ bisects $\angle BAC$, $BN\perp AN$, and $\theta$ is the measure of $\angle BAC$. If sides $AB$ and $AC$ have lengths $14$ and $19$, respectively, then find $MN$.


 
                    \begin{figure}[H]
                        centering
                        \includegraphics[width=0.5\textwidth]{Problem/AHSME_I_1981_19_Problem_Figure_1.png}
                    \end{figure}
                    


 $\textbf{(A)}\ 2\qquad\textbf{(B)}\ \dfrac{5}{2}\qquad\textbf{(C)}\ \dfrac{5}{2}-\sin\theta\qquad\textbf{(D)}\ \dfrac{5}{2}-\dfrac{1}{2}\sin\theta\qquad\textbf{(E)}\ \dfrac{5}{2}-\dfrac{1}{2}\sin\left(\dfrac{1}{2}\theta\right)$


 
 \textbf{Answer: }B



\textbf{Solution}


                    \begin{figure}[H]
                        centering
                        \includegraphics[width=0.5\textwidth]{Solution/AHSME_I_1981_19_Solution_Figure_1.png}
                    \end{figure}
                    

 Extend $BN$ to meet $AC$ at $Q$. Then $\triangle ANB \cong \triangle ANQ$. So $AQ=AB=14$, and $CQ = 19-14=5$. Also, $NB = NQ$ so $N$ is the midpoint of $\overline{BQ}$. Thus $\triangle BMN \sim \triangle BCQ$, and therefore $MN = \frac{1}{2}CQ = \boxed{(\textbf{B})\ \frac{5}{2}}$.

 


	\item (\textit{AHSME 1981}) A ray of light originates from point $A$ and travels in a plane, being reflected $n$ times between lines $AD$ and $CD$ before striking a point $B$ (which may be on $AD$ or $CD$) perpendicularly and retracing its path back to $A$ (At each point of reflection the light makes two equal angles as indicated in the adjoining figure. The figure shows the light path for $n=3$). If $\measuredangle CDA=8^\circ$, what is the largest value $n$ can have?


 
                    \begin{figure}[H]
                        centering
                        \includegraphics[width=0.5\textwidth]{Problem/AHSME_I_1981_20_Problem_Figure_1.png}
                    \end{figure}
                    


 $\textbf{(A)}\ 6\qquad\textbf{(B)}\ 10\qquad\textbf{(C)}\ 38\qquad\textbf{(D)}\ 98\qquad\textbf{(E)}\ \text{There is no largest value.}$


 
 \textbf{Answer: }B



\textbf{Solution}

Notice that when we start, we want the smallest angle possible of reflection. The ideal reflection would be $0$, but that would be impossible. Therefore we start by working backwards. Since angle $CDA$ is $8$, the reflection would give us a triangle with angles $16, 90$, and $74$. Then, when we reflect again, we will have $180 - 74 - 74$ = $32$. Since the other side of the reflection when we had the $82$ degrees had carried over to the other side, we have a $32-82-66$ triangle. 

 Notice that we keep decreasing by increments of $8$. This is because the starting angle was $8$ and since we always have to decrease $8$ every time and that every triangle has every increasing angles of $8$, we must decrease by $8$ every time. This is the most optimal path of the light beam.

 The pattern of light will be $82-74-66-58-50-42-34-26-18-10$. When we get to the angle of $2$ degrees, we have reached angle $A$. Therefore, we don't count the $2$, so our total number of reflections between $CD$ and $AD$ is $\boxed {(B) 10}$

 


	\item (\textit{AHSME 1981}) In a triangle with sides of lengths $a$, $b$, and $c$, $(a+b+c)(a+b-c) = 3ab$. The measure of the angle opposite the side length $c$ is


 $\textbf{(A)}\ 15^\circ\qquad\textbf{(B)}\ 30^\circ\qquad\textbf{(C)}\ 45^\circ\qquad\textbf{(D)}\ 60^\circ\qquad\textbf{(E)}\ 150^\circ$


 
 \textbf{Answer: }D



\textbf{Solution 1}

We will try to solve for a possible value of the variables. First notice that exchanging $a$ for $b$ in the original equation must also work. Therefore, $a=b$ works. Replacing $b$ for $a$ and expanding/simplifying in the original equation yields $4a^2-c^2=3a^2$, or $a^2=c^2$. Since $a$ and $c$ are positive, $a=c$. Therefore, we have an equilateral triangle and the angle opposite $c$ is just $\textbf{(D)}\ 60^\circ\qquad$.

 



\textbf{Solution 2}


 \[(a+b+c)(a+b-c)=3ab\]
 \[a^2+2ab+b^2-c^2=3ab\]
 \[a^2+b^2-c^2=ab\]
 \[c^2=a^2+b^2-ab\]This looks a lot like Law of Cosines, which is $c^2=a^2+b^2-2ab\cos{c}$.
 \[c^2=a^2+b^2-ab=a^2+b^2-2ab\cos{c}\]
 \[ab=2ab\cos{c}\]
 \[\frac{1}{2}=\cos{c}\]$\cos{c}$ is $\frac{1}{2}$, so the angle opposite side $c$ is $\boxed{60^\circ}$.

 -aopspandy

 


	\item (\textit{AHSME 1981}) How many lines in a three dimensional rectangular coordinate system pass through four distinct points of the form $(i,j,k)$, where $i$, $j$, and $k$ are positive integers not exceeding four?


 $\textbf{(A)}\ 60\qquad\textbf{(B)}\ 64\qquad\textbf{(C)}\ 72\qquad\textbf{(D)}\ 76\qquad\textbf{(E)}\ 100$


 
 \textbf{Answer: }D



\textbf{Solution 1}

Let $(i, j, k)$ be the first point whose coordinates are positive integers at most $4$ that a line passes through when being traced in a certain direction.  Then the next three lattice points the line passes through must be in the form $(i + a, j + b, k + c)$, $(i + 2a, j + 2b, k + 2c)$, and $(i + 3a, j + 3b, k + 3c)$, where $a, b, c$ are integers.  

 Note that if $a \geq 2$, $i + 3a \geq 1 + 3 \cdot 2 = 7$, which is too large.  Therefore $a \leq 1$, and by similar logic $b \leq 1$ and $c \leq 1$.  Also, if $a \leq -2$, $i + 3a \leq 4 + 3(-2) = -2$, which is too small.  Therefore, $a \geq -1$, and by similar logic $b \geq -1$ and $c \geq -1$.  So $a, b, c \in \{-1, 0, 1\}$.

 If $a = 1$, then $1 \leq i, i+1, i+2, i+3 \leq 4$.  In this case, only $i = 1$ is valid.

 If $a = 0$, then $1 \leq i \leq 4$.  In this case, $i = 1, 2, 3, 4$ are all valid.

 If $a = -1$, then $1 \leq i, i-1, i-2, i-3 \leq 4$.  In this case, only $i = 4$ is valid.

 Therefore, $(a, i) \in \{(1, 1), (0, 1), (0, 2), (0, 3), (0, 4), (-1, 4)\}$.  By similar logic, $(b, j)$ and $(c, k)$ must also be in this set. 

 If $a = b = c = 0$, then all four points are $(i, j, k)$, so at least one of $a, b, c$ must be nonzero.  Therefore, there are $6^3 - 4^3 = 216 - 64 = 152$ choices for $(i, j, k, a, b, c)$.  However, each line can be determined by two different values of $(i, j, k, a, b, c)$, as the line can be traced in two different directions.  For instance, $(i, j, k, a, b, c) = (4, 1, 3, -1, 1, 0)$ and $(i, j, k, a, b, c) = (1, 4, 3, 1, -1, 0)$ determine the line containing $(4, 1, 3), (3, 2, 3), (2, 3, 3), (1, 4, 3)$.  Therefore, there are $\frac{152}{2} = \boxed{(\textbf{D})\ 76}$ such lines.

 


	\item (\textit{AHSME 1981}) Equilateral $\triangle ABC$ is inscribed in a circle. A second circle is tangent internally to the circumcircle at $T$ and tangent to sides $AB$ and $AC$ at points $P$ and $Q$. If side $BC$ has length $12$, then segment $PQ$ has length


 
                    \begin{figure}[H]
                        centering
                        \includegraphics[width=0.5\textwidth]{Problem/AHSME_I_1981_23_Problem_Figure_1.png}
                    \end{figure}
                    


 $\textbf{(A)}\ 6\qquad\textbf{(B)}\ 6\sqrt{3}\qquad\textbf{(C)}\ 8\qquad\textbf{(D)}\ 8\sqrt{3}\qquad\textbf{(E)}\ 9$


 
 \textbf{Answer: }C



\textbf{Solution}

Let $O$ be the center of the smaller circle, and let $r$ be its radius.  Then $OT = OP = OQ = r$ and $AO = 2r$, since $\triangle AOP$ and $\triangle AOQ$ are $30-60-90$ triangles.  So $AT = 3r$.  Since $\triangle AOP \sim \triangle ATB$, $\frac{AP}{AB} = \frac{AO}{AT} = \frac{2}{3}$.  Since $AB = 12$, $AP = 8$ and thus $PQ = 8$. $\fbox{(C)}$.

 -j314andrews

 


	\item (\textit{AHSME 1981}) If $\theta$ is a constant such that $0 < \theta < \pi$ and $x + \dfrac{1}{x} = 2\cos{\theta}$, then for each positive integer $n$, $x^n + \dfrac{1}{x^n}$ equals


 $\textbf{(A)}\ 2\cos\theta\qquad \textbf{(B)}\ 2^n\cos\theta\qquad \textbf{(C)}\ 2\cos^n\theta\qquad \textbf{(D)}\ 2\cos n\theta\qquad \textbf{(E)}\ 2^n\cos^n\theta$


 
 \textbf{Answer: }D



\textbf{Solution}

Multiply both sides by $x$ and rearrange to $x^2-2x\cos(\theta)+1=0$. Using the quadratic equation, we can solve for $x$. After some simplifying:

 
 \[x=\cos(\theta) + \sqrt{\cos^2(\theta)-1}\]
 \[x=\cos(\theta) + \sqrt{(-1)(\sin^2(\theta))}\]
 \[x=\cos(\theta) + i\sin(\theta)\]

 Substituting this expression in to the desired $x^n + \dfrac{1}{x^n}$ gives:

 
 \[(\cos(\theta) + i\sin(\theta))^n + (\cos(\theta) + i\sin(\theta))^{-n}\]

 Using DeMoivre's Theorem:

 
 \[=\cos(n\theta) + i\sin(n\theta) + \cos(-n\theta) + i\sin(-n\theta)\]

 Because $\cos$ is even and $\sin$ is odd:

 
 \[=\cos(n\theta) + i\sin(n\theta) + \cos(n\theta) - i\sin(n\theta)\]

 
 \[=\boxed{2\cos(n\theta)},\]

 which gives the answer $\boxed{\textbf{D}}.$

 



\textbf{Solution 2}

Because we have \(x + \frac{1}{x} = 2\cos \theta\), squaring both sides gives  
 \[x^2 + 2 + \frac{1}{x^2} = 4\cos^2 \theta.\]  Subtracting 2 from both sides, we get  
 \[x^2 + \frac{1}{x^2} = 4\cos^2 \theta - 2.\]  Rewriting,  
 \[x^2 - \frac{1}{x^2} = 2(2\cos^2 \theta - 1) = 2\cos(2\theta).\]  Now, looking at the answer choices, the only one matching this form for \(n = 2\) is $\boxed{\textbf{D}}.$

 


	\item (\textit{AHSME 1981}) In $\triangle ABC$ in the adjoining figure, $AD$ and $AE$ trisect $\angle BAC$. The lengths of $BD$, $DE$ and $EC$ are $2$, $3$, and $6$, respectively. The length of the shortest side of $\triangle ABC$ is


 
                    \begin{figure}[H]
                        centering
                        \includegraphics[width=0.5\textwidth]{Problem/AHSME_I_1981_25_Problem_Figure_1.png}
                    \end{figure}
                    


 $\textbf{(A)}\ 2\sqrt{10}\qquad \textbf{(B)}\ 11\qquad \textbf{(C)}\ 6\sqrt{6}\qquad \textbf{(D)}\ 6\qquad \textbf{(E)}\ \text{not uniquely determined by the given information}$


 
 \textbf{Answer: }A



\textbf{Solution}

Let $AC=b$, $AB=c$, $AD=d$, and $AE=e$. Then, by the Angle Bisector Theorem, $\frac{c}{e}=\frac{2}{3}$ and $\frac{d}{b}=\frac12$, thus $e=\frac{3c}2$ and $d=\frac b2$.

 Also, by Stewart’s Theorem, $198+11d^2=2b^2+9c^2$ and $330+11e^2=5b^2+6c^2$. Therefore, we have the following system of equations using our substitution from earlier:

 
 \[\begin{cases}198=-\frac{3b^2}4+9c^2\\330=5b^2-\frac{75c^2}{4}\end{cases}.\]

 Thus, we have:

 
 \[\begin{cases}264=-b^2+12c^2\\264=4b^2-15c^2\end{cases}.\]

 Therefore, $5b^2=27c^2$, so $b^2=\frac{27c^2}5$, thus our first equation from earlier gives $264=\frac{33c^2}{5}$, so $c^2=40$, thus $b^2=216$. So, $c<b$ and the answer to the original problem is $c=\sqrt{40}=\boxed{2\sqrt{10}~\textbf{(A)}}$.

 



\textbf{Solution 2}

To keep the calculations simple, use a variation of stewart's formula for an angle bisector, given by: $bc=d^2+mn$.

 Using this formula and letting $AD=d$ and $AE=e$, we have $AB=2e/3$ and $AC=2d$.

 Use stewart's to obtain a system of equations, where:

 $2e^2/3 = d^2 + 6$, and

 $2d^2 = e^2 + 18$.

 Solve to get $d^2 = 54$ and $e^2 = 90$. After comparison, $2d^2 > 2e^2/3$ and we have

 $2/3 * sqrt(90) = sqrt(40) = 2sqrt(10)$

 


 


	\item (\textit{AHSME 1981}) Alice, Bob, and Carol repeatedly take turns tossing a die. Alice begins; Bob always follows Alice; Carol always follows Bob; and Alice always follows Carol. Find the probability that Carol will be the first one to toss a six. (The probability of obtaining a six on any toss is $\dfrac{1}{6}$, independent of the outcome of any other toss.)


 $\textbf{(A)}\ \dfrac{1}{3}\qquad \textbf{(B)}\ \dfrac{2}{9}\qquad \textbf{(C)}\ \dfrac{5}{18}\qquad \textbf{(D)}\ \dfrac{25}{91}\qquad \textbf{(E)}\ \dfrac{36}{91}$


 
 \textbf{Answer: }D



\textbf{Solution 1 (Infinite Series)}

For Carol to win on her $n^{th}$ roll, the first $3n-1$ rolls must not be $6$ and the last roll must be $6$.  Therefore, the probability that Carol wins on her $n^{th}$ roll is $\frac{1}{6} \cdot \left(\frac{5}{6}\right)^{3n-1} = \frac{1}{5} \cdot \left(\frac{125}{216}\right)^{n}$.  Taking the sum of this expression over all positive integers $n$, her total probability of winning must be $\sum_{n=1}^{\infty} \frac{1}{5} \cdot \left(\frac{125}{216}\right)^{n}$.  Since this is a geometric series with first term $\frac{25}{216}$ and common ratio $\frac{125}{216}$, Carol's probability of winning is  $\frac{\frac{25}{216}}{1-\frac{125}{216}}=\frac{\frac{25}{216}}{\frac{91}{216}}=\frac{25}{91}$. $\fbox{(D)}$

 



\textbf{Solution 2 (Probability States)}

Let $a$, $b$, and $c$ be Carol's probabilities of winning, depending on whether it is Alice's, Bob's or Carol's turn to roll the die next, respectively.  Since it is initially Alice's turn, the value of $a$ will be the answer.

 Then $a = \frac{5}{6}b$, $b = \frac{5}{6}c$, and $c = \frac{1}{6} + \frac{5}{6}a$.  Substituting into the first equation shows that $a = \frac{5}{6}b = \frac{25}{36}c = \frac{25}{36}\left(\frac{1}{6} + \frac{5}{6}a\right) = \frac{25}{216} + \frac{125}{216}a$. Solving the equation $a = \frac{25}{216} + \frac{125}{216}a$ gives $a = \frac{25}{91}$. $\fbox{(D)}$ 

 -j314andrews

 


	\item (\textit{AHSME 1981}) In the adjoining figure triangle $ABC$ is inscribed in a circle. Point $D$ lies on $\stackrel{\frown}{AC}$ with $\stackrel{\frown}{DC} = 30^\circ$, and point $G$ lies on $\stackrel{\frown}{BA}$ with $\stackrel{\frown}{BG}\, > \, \stackrel{\frown}{GA}$. Side $AB$ and side $AC$ each have length equal to the length of chord $DG$, and $\angle CAB = 30^\circ$. Chord $DG$ intersects sides $AC$ and $AB$ at $E$ and $F$, respectively. The ratio of the area of $\triangle AFE$ to the area of $\triangle ABC$ is


 
                    \begin{figure}[H]
                        centering
                        \includegraphics[width=0.5\textwidth]{Problem/AHSME_I_1981_27_Problem_Figure_1.png}
                    \end{figure}
                    


 $\textbf{(A)}\ \dfrac {2 - \sqrt {3}}{3}\qquad \textbf{(B)}\ \dfrac {2\sqrt {3} - 3}{3}\qquad \textbf{(C)}\ 7\sqrt {3}-12\qquad \textbf{(D)}\ 3\sqrt {3}-5\qquad\\ \textbf{(E)}\ \dfrac {9-5\sqrt {3}}{3}$


 
 \textbf{Answer: }C



\textbf{Solution}

Since $\triangle ABC$ is isosceles, $\angle ACB = \angle ABC = \frac{150^{\circ}}{2} = 75^{\circ}$, and $\stackrel{\frown}{AC}\ =\ \stackrel{\frown}{AB}\ = 2 \cdot 75^{\circ} = 150^{\circ}$.  Thus $\stackrel{\frown}{AD}\ = 150^{\circ} - 30^{\circ} = 120^{\circ}$, and $\angle AGD = \angle ACD = \frac{120^{\circ}}{2} = 60^{\circ}$.  Since $DG = AB = AC$, $\stackrel{\frown}{DG}\ = 150^{\circ}$.  So $\stackrel{\frown}{AG}\ = 150^{\circ} - 120^{\circ} = 30^{\circ}$.  Therefore, $\angle AEG = \angle DEC = \frac{30^{\circ} + 30^{\circ}}{2} = 30^{\circ}$, and $\triangle AFE$ is isosceles with $AF = FE$.  Also $\angle AFE = 180^\circ - 2 \cdot 30^\circ = 120^\circ$ and $\angle AFG = 60^\circ$, so $\triangle AFG$ is equilateral.

 Let $x = AF = FE = FG$.  By the Law of Cosines, $AE^2 = x^2 + x^2 - 2x \cdot x \cdot \cos(120^\circ) = 3x^2$, so $AE = x\sqrt{3}$.  $\stackrel{\frown}{AG}\ =\ \stackrel{\frown}{DC}$, $AG = DC$.  Thus $\triangle DEC \cong \triangle AEG$, and $EC = 2x$.  Thus $AC = AB = x(2+\sqrt{3})$.  The area of $\triangle AFE$ is $\frac{1}{2} \cdot x \cdot x\sqrt{3} \cdot \sin 30^\circ = \frac{x^2\sqrt{3}}{4}$, while the area of $\triangle ABC$ is $\frac{1}{2} \cdot x(2+\sqrt{3}) \cdot x(2+\sqrt{3}) \cdot \sin 30^\circ = \frac{x^2(7+4\sqrt{3})}{4}$.  

 The ratio of these two areas is $\frac{\frac{x^2\sqrt{3}}{4}}{\frac{x^2(7+4\sqrt{3})}{4}} = \frac{\sqrt{3}}{7+4\sqrt{3}} = \sqrt{3}(7-4\sqrt{3}) = 7\sqrt{3} - 12$ $\fbox{(D)}$.

 -j314andrews

 


	\item (\textit{AHSME 1981}) Consider the set of all equations $x^3 + a_2x^2 + a_1x + a_0 = 0$, where $a_2$, $a_1$, $a_0$ are real constants and $|a_i| \leq 2$ for $i = 0,1,2$. Let $r$ be the largest positive real number which satisfies at least one of these equations. Then


 $\textbf{(A)}\ 1 \leq r < \dfrac{3}{2}\qquad \textbf{(B)}\ \dfrac{3}{2} \leq r < 2\qquad \textbf{(C)}\ 2 \leq r < \dfrac{5}{2}\qquad \textbf{(D)}\ \dfrac{5}{2} \leq r < 3\qquad \textbf{(E)}\ 3 \leq r < \dfrac{7}{2}$


 
 \textbf{Answer: }D



\textbf{Solution}

Since $x^3 = -(a_2x^2 + a_1x + a_0)$ and $x$ will be as big as possible, we need $x^3$ to be as big as possible, which means $a_2x^2 + a_1x + a_0$ is as small as possible. Since $x$ is positive (according to the options), it makes sense for all of the coefficients to be $-2$.

 Evaluating $f(\frac{5}{2})$ gives a negative number, $f(3)$ 1, and $f(\frac{7}{2})$ a number greater than 1, so the answer is $\boxed{D}$

 


	\item (\textit{AHSME 1981}) If $a \geq 1$, then the sum of the real solutions of 


 $\sqrt{a - \sqrt{a + x}} = x$


 is equal to 


 $\textbf{(A)}\ \sqrt{a} - 1\qquad \textbf{(B)}\ \dfrac{\sqrt{a}- 1}{2}\qquad \textbf{(C)}\ \sqrt{a - 1}\qquad \textbf{(D)}\ \dfrac{\sqrt{a - 1}}{2}\qquad \textbf{(E)}\ \dfrac{\sqrt{4a- 3} - 1}{2}$


 
 \textbf{Answer: }E



\textbf{Solution 1}

A solution is available \href{https://files.eric.ed.gov/fulltext/ED239856.pdf}{here}. Pull up find, and put in "Since x is the principal", and you will arrive at the solution.

 It's not super clear, and there's some black stuff over it, but its legible. 

 The solution in  the above file/pdf is the following. I tried my best to match it verbatim, but I had to guess at some things. I also did not do the entire solution like this, just parts where I had to figure out what the words/math was, so this transcribed solution could be wrong and different from the solution in the aforementioned file/pdf.

 Anyways:

 29. (E) Since $x$ is the principal square root of some quantity, $x\geq0$. For $x\geq0$, the given equation is equivalent to 
 \[a-\sqrt{a+x}=x^2\] or 
 \[a=\sqrt{a+x}+x^2.\] The left member is a constant, the right member is an increasing function of $x$, and hence the equation has exactly one solution. We write 
 \begin{align*} \sqrt{a+x}&=a-x^2 \\ \sqrt{a+x}+x&=(a+x)-x^2 \\ &=(\sqrt{a+x}+x)(\sqrt{a+x}-x). \end{align*}


 Since $\sqrt{a+x}+x>0$, we may divide by it to obtain 
 \[1=\sqrt{a+x}-x\quad\text{or}\quad x+1=\sqrt{a+x},\]so 
 \[x^2+2x+1=a+x,\] and 
 \[x^2+x+1-a=0.\]

 Therefore $x=\frac{-1\pm\sqrt{4a-3}}{2}$, and the positive root is $x=\frac{-1+\sqrt{4a-3}}{2}$, the only solution of the original equation. Therefore, this is also the sum of the real solutions.

 -OliverA

 



\textbf{Solution 2}

Recall that the square root of a real number cannot be negative, so $x \geq 0$.  Square both sides of the equation to get $a - \sqrt{a + x} = x^2$, which is equivalent to $a - x^2 = \sqrt{a+x}$.    

 Once again, the square root of a real number cannot be negative, so $a - x^2 \geq 0$, that is $x^2 \geq a$.  So $-\sqrt{a} \leq x \leq \sqrt{a}$, and since $x \geq 0$ as previously determined, $0 \leq x \leq \sqrt{a}$.  

 Squaring both sides again yields $a^2 - 2ax^2 + x^4 = a + x$.  Collecting all terms on the yields $a^2 - 2ax^2 - a + x^4 - x =$ $a^2 - (2x^2 + 1) a + (x^4 - x) = 0$.  

 Using the quadratic formula, $a = \frac{2x^2+1 \pm \sqrt{(2x^2+1)^2 - 4(x^4 -x)}}{2} =$ $\frac{2x^2 + 1 \pm \sqrt{4x^2 + 4x + 1}}{2} = \frac{2x^2 + 1 \pm (2x + 1)}{2}$, so either $a = x^2 + x + 1$ or $a = x^2 - x$.

 If $a = x^2 + x + 1$, then $x^2 + x + (1-a) = 0$, and by the quadratic formula, $x = \frac{-1 \pm \sqrt{1 - 4(1-a)}}{2} = \frac{-1 \pm \sqrt{4a - 3}}{2}$.  Since $\frac{-1-\sqrt{4a-3}}{2}$ is negative, it cannot a valid solution.  However, $\frac{-1+\sqrt{4a-3}}{2} \geq \frac{-1+1}{2} = 0$, and $\frac{-1+\sqrt{4a-3}}{2} < \frac{-1 + 2\sqrt{a}}{2} = -\frac{1}{2} + \sqrt{a} < \sqrt{a}$, so $\frac{-1+\sqrt{4a-3}}{2}$ is a valid solution.

 If $a = x^2 - x$, then $x^2 - x - a = 0$, and by the quadratic formula, $x = \frac{1 \pm \sqrt{1+4a}}{2}$.  But $\frac{1-\sqrt{1+4a}}{2} < \frac{1-1}{2} = 0$, and $\frac{1+\sqrt{1+4a}}{2} > \frac{\sqrt{1+4a}}{2} > \frac{\sqrt{4a}}{2} = \sqrt{a}$, so neither of these are valid solutions.

 Therefore, $x = \frac{-1+\sqrt{4a-3}}{2}$ is the only solution to the original equation, so the sum of all solutions is $\boxed{(\textbf{E})\ \frac{-1+\sqrt{4a-3}}{2}}$.

 -OliverA, edited by j314andrews.  Similar to the alternate solution from The Contest Problem Book IV.

 \textit{Note}: One might notice that when $a=3$, the solution of the original equation is $x=1$. This eliminates all choices except (E).

 -- OliverA

 


	\item (\textit{AHSME 1981}) If $a$, $b$, $c$, and $d$ are the solutions of the equation $x^4 - bx - 3 = 0$, then an equation whose solutions are 

 \[\dfrac {a + b + c}{d^2}, \dfrac {a + b + d}{c^2}, \dfrac {a + c + d}{b^2}, \dfrac {b + c + d}{a^2}\]is


 $\textbf{(A)}\ 3x^4 + bx + 1 = 0\qquad \textbf{(B)}\ 3x^4 - bx + 1 = 0\qquad \textbf{(C)}\ 3x^4 + bx^3 - 1 = 0\qquad \\\textbf{(D)}\ 3x^4 - bx^3 - 1 = 0\qquad  \textbf{(E)}\ \text{none of these}$


 
 \textbf{Answer: }D



\textbf{Solution 1}

Using Vieta's formula, we know the sum of the roots is equal to the negative coefficient of the $x^3$ term. Since the coefficient is 0, $a+b+c+d=0$. Thus, $\frac{a+b+c}{d^2}$ can be rewritten as $\frac{-d}{d^2}=\frac{1}{-d}$. Similarly, the other three new roots can be written as $\frac{1}{-c}$, $\frac{1}{-b}$, and $\frac{1}{-a}$.

 Now, we need to find a way to transform the function $f(x)=x^4-bx-3$ such that all the roots are its negative reciprocal. We can create this new function by taking the negative reciprocal of the argument. In other words, $f(\frac{1}{-x})$ satisfies this criteria.

 The new equation, $f(\frac{1}{-x})=0$ has the required roots and can be simplified to $\frac{1}{x^4}+\frac{b}{x}-3=0$. Since this is not a polynomial, we can multiply both sides by $x^4$ to become $1+bx^3-3x^4=0$. After rearranging and multiplying by negative one, we arrive at $3x^4-bx^3-1$ so the answer is $\boxed{\textbf{(D)} 3x^4-bx^3-1}$

 



\textbf{Solution 2}

As in solution 1, the roots of the new equation are $-\frac{1}{a}, -\frac{1}{b}, -\frac{1}{c},-\frac{1}{d}$. Furthermore, applying Vieta’s formula to the original equation yields $abcd=-3$ and $abc+abd+acd+bcd=-3(\frac{1}{a}+\frac{1}{b}+\frac{1}{c}+\frac{1}{d})=b$. Therefore, the product of the zeros of the new equation is $\frac{1}{abcd}=-\frac{1}{3}$. This limits our choices to options C and D, and we need to look for the sum of the roots of the new equation. This sum equals to $-(\frac{1}{a}+\frac{1}{b}+\frac{1}{c}+\frac{1}{d})$The function whose roots are the reciprocals of the original equation is $-3x^4-bx^3+1$ therefore $-(\frac{1}{a}+\frac{1}{b}+\frac{1}{c}+\frac{1}{d})=-(-\frac{-b}{-3})=\frac{b}{3}$. The second term of the chosen equation $a_{2}$ should satisfy that $-\frac{a_{2}}{3}=\frac{b}{3}$, hence $a_{2}=-b$. The answer is $D$. (Option E looks ridiculous~)

 


\end{enumerate}